%% ════════════════════════════════════════════════════════════════════════════
\section{Conclusion}
\label{sec:conclusion}

This paper presented an empirical audit of data leakage in Indian air quality forecasting literature and established, to our knowledge, the first leak-free, rolling-origin benchmark across Gujarat cities. We demonstrated that the combination of concurrent pollutant inputs, random train-test splitting, and temporal regime shifts can produce severe apparent performance inflation -- exemplified by Random Forest in Ahmedabad, where explained variance ($R^2$) fell from $0.815$ under Setup~A to $0.015$ under Setup~B (though excluding the March--May 2020 COVID-19 lockdown days from Setup~B restores $R^2$ to $0.747$, showing that distribution shift contributed substantially alongside autocorrelation leakage).

Under rigorous, leak-free rolling-origin evaluation at next-day lead times ($h=1$), machine learning models failed to statistically outperform simple classical baselines (with Persistence the simplest, and ARIMA(1,0,0) matching or beating ML in Gandhinagar and Surat) in five of six evaluated cities. In industrial Vapi, Persistence was itself statistically superior to Random Forest and XGBoost ($p \leq 0.003$). At $h=1$, machine learning established nominally significant skill only in coastal Chennai ($+11.5\%$ skill, $p = 0.002$), though this does not survive Bonferroni correction under $k=0$ ($\alpha = 0.05/60$; XGBoost is borderline at $p = 0.00075$ under $k=3$). Genuine Bonferroni-corrected skill is confirmed at multi-day lead times ($h=3$) in Chennai (Ridge $p = 0.0006$, XGBoost $p = 0.0002$) where atmospheric autocorrelation naturally attenuates.

These findings demonstrate that future research claiming algorithmic superiority for urban air quality forecasting should benchmark against persistence under strict temporal validation. To support community adoption, our complete pipeline, datasets, and evaluation ledgers are openly accessible.

%% ════════════════════════════════════════════════════════════════════════════
\section*{Acknowledgements}
The author acknowledges the Central Pollution Control Board (CPCB), Ministry of Environment, Forest and Climate Change, Government of India, for maintaining the public air quality monitoring network and making municipal bulletin records publicly accessible.

\section*{Reproducibility Statement}
The entire experimental pipeline is designed for full, single-command reproduction. Source code, configuration files, cleaned datasets, and pre-computed prediction ledgers are openly accessible at: \url{https://github.com/Pratham1227-beep/aqi-gujarat-leakage}. Complete execution commands, dependencies, and reproduction instructions are provided in the repository README.

%% ════════════════════════════════════════════════════════════════════════════
